\begin{statement}{P3K.59}{Refuted; physical-line correction}\par\noindent\textbf{Physical-line correction, 17 November 2026.} The actual two central sequences are \(-1-3j(j+1)\) and \(-7-24j(j+1)\). The second has least period seven and image \(\{1,13,17,21\}\). All displayed matrices, native free pairs and directed mutation identities are retained. The old additive baselines are refuted and corrected in K69. The following earlier assertion and its argument are retained as the located historical claim; their corrected current consequences are proved in K68--69.\medskip\par 
Retain the actual positive selected-cusp source loop of
P3F.54 and P3F.56, with its fixed scalar, prescribed finite cap,
physical height, Cartier line and complete period units. For every
integer \(j\geq0\), consider the two specified periodic fans with bases
\[
 P_j=\begin{pmatrix}1&-j\\0&1\end{pmatrix},\qquad
 R_j=\begin{pmatrix}j+1&-j\\-1&1\end{pmatrix}.
\]
Their actual oriented clutched-sphere classes are
\[
 [\Phi_{P_j}]=1-3j(j+1),\qquad
 [\Phi_{R_j}]=-5-24j(j+1)\quad\text{in }\Theta_7=\mathbb Z/28.
\]
On \(P_j\), the complete ordinary-product variation and ambient
stable-derivative invariants in the marked three-curve basis are
\[
 V_j=\begin{pmatrix}
 0&6(j+1)&-6\\
 6(j+1)&18j(j+1)&-6j\\
 -6&-6j&0
 \end{pmatrix},\qquad
 \delta_j=\begin{pmatrix}-12(j+1)\\-24j(j+1)-12\\12j\end{pmatrix}.
\]
The corresponding native relative pairs, in the order \(C_0,C_1,C_2\), are
\[
 (6j,-6(j+1)),\qquad
 (3j(j+1),-21j(j+1)-6),\qquad
 (-6j-6,6j).
\]
For each selected periodic flop used in these constructions, write
\(a=V_{ii}\), \(b=\delta_i\), and let \(i'\) be the new label of that
same flopped curve. The actual native Cox arrow and the carried
ordinary-product comparison give exactly
\[
 V'_{i'i'}=a,\qquad \delta'_{i'}=2a-b,\qquad
 s'=s+\frac{b-a}{2}\pmod {28}.
\]
The \(R_j\) return sequence has least positive period seven and exact
range \(\{3,15,19,23\}\). In particular \(R_1\) has class three.
The statement concerns these explicitly constructed families and their
selected updates; it assigns no return to an additional fan.
\end{statement}
\begin{proof}
\par\medskip\noindent\textbf{CF.1. Three exact lattice and native-star arrows.}\par\nopagebreak[4]

Write the ordered fan basis as \(P=(v\ w)\). Its triangles are
\((0,v,w)\) and \((v,w,v+w)\), with their lattice translates.
In the first triangle use \((a,b,c)\), with
\(q=abc, Z=b,V=c\). The three curves have native first charts
\[
 C_0:(z,n_0,n_1)=(b,a,c),\qquad
 C_1:(z,n_0,n_1)=(c,a,b),\qquad
 C_2:(z,n_0,n_1)=(a,b,c).                              \tag{P3K.59.CF.1}
\]
Their native second charts are exactly GF.11 and satisfy
\(z'=z^{-1},n'_r=zn_r\), including every complete period unit.

The three periodic flips have the following actual new bases and
new labels for the flopped curve:
\[
 Q_0=\begin{pmatrix}1&-1\\0&1\end{pmatrix},\qquad
 Q_1=\begin{pmatrix}1&0\\-1&1\end{pmatrix},\qquad
 Q_2=\begin{pmatrix}1&1\\0&1\end{pmatrix},
 \qquad P'=PQ_i,\qquad i'=(2,2,0)_i.                  \tag{P3K.59.CF.2}
\]
For \(C_0\), the old quadrilateral is
\((0,v,w,w-v)\); its old diagonal is \((0,w)\), and its new
diagonal joins \(v\) to \(w-v\). The new two triangles are those
of \((v,w-v)\). For \(C_1\), use
\((0,w,v,v-w)\), giving the basis \((v-w,w)\).
For \(C_2\), the quadrilateral is \((0,v,w,v+w)\), and its
new diagonal is \((0,v+w)\). The second triangle of
\((v,v+w)\), translated by \(-v\), is \((0,w,v+w)\).
These are literal triangle equalities, repeated over the period
lattice. Each basis change and its inverse is integral.

The actual torus characters are, respectively,
\[
 (Z',V')=(ZV,V),\qquad (Z',V')=(Z,ZV),\qquad
 (Z',V')=(Z/V,V),\qquad q'=q.                         \tag{P3K.59.CF.3}
\]
Their new first-triangle coordinates are
\[
 (a/c,bc,c),\qquad (a/b,b,bc),\qquad (ac,b/c,c).        \tag{P3K.59.CF.4}
\]
Thus in the native first charts of the flopped curves, (P3K.59.CF.3)
is exactly the standard other-ruling arrow
\[
 (z,n_0,n_1)\longmapsto
  (\zeta,u,v)=(n_0/n_1,zn_1,n_1)                     \tag{P3K.59.CF.5}
\]
for \(i=0,2\). For \(i=1\) the last two coordinates are
\((v,u)\). In each case \(\zeta uv=zn_0n_1=q\).
The second new chart uses
\((\zeta',u',v')=(n_1/n_0,zn_0,n_0)\), with the same
constant interchange for \(i=1\). It has
\(\zeta'=\zeta^{-1},(u',v')=\zeta(u,v)\).
These formulas are the two quotient charts of the rank-one matrix
\[
 \begin{pmatrix}n_0&n_1\\zn_0&zn_1\end{pmatrix}.       \tag{P3K.59.CF.6}
\]
They identify the common complement and both native normal bundles
by the actual Cox construction. For the normal interchange at
\(i=1\), use the explicit path
\[
 S_t=\frac12
 \begin{pmatrix}
 1+e^{\pi it}&1-e^{\pi it}\\
 1-e^{\pi it}&1+e^{\pi it}
 \end{pmatrix},\qquad 0\le t\le1.                    \tag{P3K.59.CF.6a}
\]
Its eigenspaces are the two constant lines spanned by
\((1,1)\) and \((1,-1)\), with eigenvalues \(1,e^{\pi it}\).
Thus \(S_t^*S_t=1\), \(S_0=1\), and
\(S_1=(\begin{smallmatrix}0&1\\1&0\end{smallmatrix})\).
It commutes with the transition \(z\,1\) and therefore acts on
the whole native normal bundle. It is independent of the source
loop parameter. Its oriented real lift conjugates the two
Spin factors through constant conjugations, which preserve
both integral third degrees.

Here are the full unit maps. Write the old two periods as GF.2,
with unit matrix \(A=(A_{rs})\). The new generators are the
columns of \(Q_i\) in that period basis. The actual new units are
\[
 A'_{rs}=\prod_{h,j}A_{hj}^{(Q_i^{-1})_{rh}(Q_i)_{js}}. \tag{P3K.59.CF.7}
\]
In the three cases their matrices are
\[
 A'_0=\begin{pmatrix}
 A_{11}A_{21}&A_{12}A_{22}/(A_{11}A_{21})\\
 A_{21}&A_{22}/A_{21}
 \end{pmatrix},
\]
\[
 A'_1=\begin{pmatrix}
 A_{11}/A_{12}&A_{12}\\
 A_{11}A_{21}/(A_{12}A_{22})&A_{12}A_{22}
 \end{pmatrix},
\]
\[
 A'_2=\begin{pmatrix}
 A_{11}/A_{21}&A_{11}A_{12}/(A_{21}A_{22})\\
 A_{21}&A_{21}A_{22}
 \end{pmatrix}.                                      \tag{P3K.59.CF.8}
\]
Substitution in the new GF.11 second charts gives the second
native chart of (P3K.59.CF.5) exactly. For example in case \(i=2\),
the second vertex coordinates before applying GF.11 are
\[
 (u',v',w')=\bigl(A_{11}b,ab/A_{21},A_{21}c/(A_{11}b)\bigr).
\]
GF.11 sends these to \((c/b,ab,b)\), precisely the second
chart of (P3K.59.CF.5). For \(i=0,1\), the new curve is \(C_2\),
and the second native charts are directly
\((c/a,ab,a)\) and \((b/a,a,ac)\), respectively.
These equalities retain all unit factors; their parameter winding
is included in (P3K.59.CF.7). In particular
\(B'=Q_i^{-1}BQ_i\), without replacing a unit by its winding.

The characters (P3K.59.CF.3), inverse characters and full period changes
give the same original source and line maps on their dense overlap.
The physical height is literally \(q\). The retained finite cap
and source exterior therefore use their original maps there.
Every periodic curve orbit has the same free representative charts
as K55, so (P3K.59.CF.5)--\allowbreak (P3K.59.CF.8) descend to its whole periodic flop.
The unrestricted four-weight Cox identity of K54 intertwines
the physical circle with these maps. Transporting the Hermitian
metric and applying the positive radial normalization gives its
actual common-link arrow; the normal rescaling of K54 carries
the full nonlinear units. This supplies the actual family flop,
including its native normal conventions.

\par\medskip\noindent\textbf{CF.2. Exact selected-coordinate update.}\par\nopagebreak[4]

Let the old ordinary relative pair at \(C_i\) be \((m_i,n_i)\),
and let its central coordinate be the actual sphere return \(s\).
K55's explicit cylinder comparison removes the pure-left
\(L^6\), giving the native pair \((m_i-6,n_i)\).
The actual common-link arrow (P3K.59.CF.5), in the K23 Hopf coordinates,
is \(\tau(h,q_H)=(q_H^{-1}hq_H,q_H)\), with the exterior
extension \(T(y,q_H)=(q_H^{-1}yq_H,q_H)\). The complex-normal
reset is exactly K54's fixed \(J\). Applying their proved
relative maps gives the new native pair and central coordinate
\[
 (m_i-6+2n_i,-n_i),\qquad s'=s+n_i.                   \tag{P3K.59.CF.9}
\]
The constant normal interchange in (P3K.59.CF.5) acts by a constant
oriented four-plane change. Its lift conjugates the two Spin
factors separately, through the displayed \(U(2)\) path, and
preserves their degrees. Thus the same formula applies to
all three cases of (P3K.59.CF.2).

Add the actual new \(L^6\) comparison on the new native star.
Its ordinary pair is \((m_i+2n_i,-n_i)\). In terms of the
actual ordinary variation and ambient derivative, put
\(a=V_{ii}\), \(b=\delta_i\). K31.7's carried single-star
frame map gives
\[
 a=-(m_i+n_i),\qquad b=2n_i+a.
\]
Consequently the new selected entries and the actual return obey
\[
 V'_{i'i'}=a,\qquad \delta'_{i'}=2a-b,\qquad
 s'=s+\frac{b-a}{2}\pmod {28}.                       \tag{P3K.59.CF.10}
\]
This derivation uses only the actual selected native star on
each side. It does not use a simultaneous-product matrix reset.

\par\medskip\noindent\textbf{CF.3. Recovering the other entries from the actual trace maps.}\par\nopagebreak[4]

For every target with nonzero tangent weights \(w_0,w_1,w_2\),
write \(r_i=1/w_i\). The exact GF trace identities and symmetry
have the complete solutions
\[
 V=\operatorname{diag}(I_0/w_0,I_1/w_1,I_2/w_2)+c\,rr^T,
 \qquad \delta_j=(W_j+e)/w_j.                         \tag{P3K.59.CF.11}
\]
Indeed the displayed diagonal is a particular solution. The
two vectors \(w_je_j-w_0e_0\) span the kernel of the functional
with coefficients \(r\). A symmetric matrix annihilating that
plane has image in its orthogonal line, and is exactly
\(c rr^T\). The same argument for the derivative leaves
the single constant \(e\). All these statements are rational
linear algebra applied to the actual integral identities.

At the selected new curve the already proved entries (P3K.59.CF.10)
therefore fix
\[
 c=w_{i'}^2a-w_{i'}I_{i'},\qquad
 e=w_{i'}(2a-b)-W_{i'}.                               \tag{P3K.59.CF.12}
\]
Equations (P3K.59.CF.11)--\allowbreak (P3K.59.CF.12) give the entire actual new matrix
and derivative. Their values are integral because they equal
the actual integral invariants uniquely, whose existence is
the explicit family construction above.

For completeness a target has at most one zero tangent weight.
For the zero index \(r\), the other two weights are \(+1,-1\),
by primitivity. The same trace equations give
\[
 \begin{aligned}
 V_{jj}&=I_j/w_j,\qquad V_{rj}=V_{jr}=-I_r/w_j,\\
 V_{pq}&=0\qquad(p,q\ne r,\ p\ne q),\\
 \delta_j&=(W_j-W_r)/w_j\qquad(j\ne r).
 \end{aligned}                                      \tag{P3K.59.CF.13}
\]
Their two remaining entries are precisely \(V_{rr},\delta_r\).
For a flop whose selected target is that fixed curve, (P3K.59.CF.10)
supplies both. A target with a different retained fixed curve
requires its own actual fixed-star data. This describes the
exact rank of the trace map and leaves that specific data open.

\par\medskip\noindent\textbf{CF.4. A complete horizontal staircase.}\par\nopagebreak[4]

For every integer \(j\ge0\), set
\[
 P_j=\begin{pmatrix}1&-j\\0&1\end{pmatrix}.             \tag{P3K.59.CF.14}
\]
The actual return and complete ordinary invariants on this fan are
\[
 s_j=1-3j(j+1)\pmod {28},
\]
\[
 V_j=\begin{pmatrix}
 0&6(j+1)&-6\\
 6(j+1)&18j(j+1)&-6j\\
 -6&-6j&0
 \end{pmatrix},\qquad
 \delta_j=\begin{pmatrix}
 -12(j+1)\\-24j(j+1)-12\\12j
 \end{pmatrix}.                                      \tag{P3K.59.CF.15}
\]
The native pairs on its three curves are
\[
 (6j,-6(j+1)),\quad
 (3j(j+1),-21j(j+1)-6),\quad
 (-6j-6,6j).                                         \tag{P3K.59.CF.16}
\]

Here is an induction by actual family maps. At \(j=0\), (P3K.59.CF.15)
is exactly K51/K52 and the return is F56's class one. GF.10a
supplies the explicit based homotopy identifying those initial
products with K55's chosen product. The \(C_0\) flip of \(P_j\)
is precisely \(P_jQ_0=P_{j+1}\); its selected target is \(C_2\).
For \(P_j\) the literal complete windings and trace degrees are
\[
 B_j=\begin{pmatrix}-6&6j\\0&0\end{pmatrix},\quad
 l=(0,-6(j+1),0),\quad w=(j,1,-j-1),
\]
\[
 I=(-6(j+1),12j(j+1),6j),\qquad
 W=(6,-12j(j+1)-6,6).                                \tag{P3K.59.CF.17}
\]
These follow by substituting the actual (P3K.59.CF.3)/(P3K.59.CF.7) maps in GF.
The target weights \((j+1,1,-j-2)\) are all nonzero.
Equation (P3K.59.CF.10) sets its selected entries to
\(V'_{22}=0,\delta'_2=12(j+1)\). The unique solution
(P3K.59.CF.11)--\allowbreak (P3K.59.CF.12) is (P3K.59.CF.15) with \(j\) replaced by \(j+1\).
Thus the full matrix and derivative have been obtained from the
actual target, including their off-diagonal entries.
The old right integer is \(-6(j+1)\), so
\[
 s_{j+1}=s_j-6(j+1).
\]
Summing this actual consecutive-flop identity gives (P3K.59.CF.15).
Finally the native comparison gives
\(n_i=(\delta_i-V_{ii})/2\) and
\(m_i=-(\delta_i+V_{ii})/2-6\), which proves (P3K.59.CF.16).

\par\medskip\noindent\textbf{CF.5. The adjacent complete family.}\par\nopagebreak[4]

Flop \(C_1\) on each of the just constructed actual families.
Its new basis is
\[
 R_j=P_jQ_1=\begin{pmatrix}j+1&-j\\-1&1\end{pmatrix},
 \qquad j\ge0.                                       \tag{P3K.59.CF.18}
\]
Equation (P3K.59.CF.16) gives its actual native right integer.
The same Cox family arrow (P3K.59.CF.5)--\allowbreak (P3K.59.CF.10) therefore gives
\[
 [\Phi_{R_j}]=1-3j(j+1)-21j(j+1)-6
             =-5-24j(j+1)\pmod {28}.                 \tag{P3K.59.CF.19}
\]
This computes the actual return on every fan (P3K.59.CF.18), including
\(j=0\). Its sequence has period seven: modulo 28 it equals
\(23+4j(j+1)\), and shifting \(j\) by seven changes this
by \(56j+224\). The seven consecutive values are
\[
                         23,3,19,15,19,3,23 .
\]
They give the exact range \(\{3,15,19,23\}\); the sequence
is not constant, so seven is its least positive period.
In particular
\(R_1=(\begin{smallmatrix}2&-1\\-1&1\end{smallmatrix})\)
has the actual positive return class three.

All original physical source and cap maps are carried through
the specified finite sequence for each \(j\). This proof assigns
no value to a fan outside its computed domain. The exact algebra
and repeated-basis/fan consistency checks are retained under
\path{code/part3_root/source_cusp_mutation_diagnostic.*}; their
first unsimplified symbolic-matrix equality attempt is preserved.
The written Cox/period maps, native-frame comparison and induction
are the proof for all \(j\). The active finite search starts only
from the proved \(P_0=1\) seed: at depth seven it checks
210 basis states and 128 fan states, with no conflicting state
or failed integrality check. Four transitions to a target with
a separately retained fixed curve are left unevaluated.
The earlier exploratory two-seed output is retained separately;
its second fixed-star seed supplies no input to this proof.
\end{proof}
